%\setlength{\columnsep}{1cm}
%\setlength{\parskip}{1ex plus 0.5ex}
%para editar el tamaño de los margenes
\usepackage[a4paper,top=2.54cm,bottom=2.54cm,left=2.54cm,right=2.54cm,marginparwidth=1.75cm]{geometry}
%\setlength{\parindent}{.5 cm}   %Activar sangrias
\usepackage{latexsym,url,colortbl,amsmath,amssymb,amsfonts,amsthm,mdframed,fancyhdr,lastpage}
\usepackage{paralist,etoolbox,enumerate,mathrsfs}
\usepackage[usenames,dvipsnames]{xcolor}
\pagestyle{fancy}\fancyhf{}
%\renewcommand{\headrulewidth}{0pt}}
%\spanishdecimal{.}
\pretolerance=2000
\tolerance=3000
\usepackage[pdftex]{graphicx}
\usepackage{csquotes}
\counterwithin*{equation}{section}
\counterwithin*{equation}{subsection}
\renewcommand\theequation{\ifnumgreater{\value{subsection}}{0}{\thesubsection.}{\thesection.}\arabic{equation}}%

\usepackage[backend=biber, style=alphabetic]{biblatex}
\addbibresource{referencias.bib}
\usepackage{colortbl}
\usepackage{booktabs}
%\usepackage{xcolor}
\usepackage{array, multirow, multicol}
\usepackage{color}

\usepackage{float}
\usepackage{subcaption}
\usepackage{amsmath,amssymb,amsthm} % Paquetes de matemáticas
\usepackage{amsfonts}           % Fuentes matemáticas adicionales
\usepackage{geometry}           % Control del formato de página
\usepackage{hyperref}           % Enlaces y referencias internas
\usepackage{mathrsfs}           % Fuente caligráfica
\usepackage{wasysym}
\usepackage{fancyhdr}           % Encabezados y pies de página personalizados
\usepackage{tikz-cd}
\usepackage{tikz}
\usetikzlibrary{babel}
\usepackage{
abc}
\usepackage[utf8]{inputenc}
\usepackage[T1]{fontenc}
\usepackage{helvet}

\usepackage[hidelinks]{hyperref}
%\renewcommand{\familydefault}{\sfdefault}

%\usepackage[backend=bibtex,style=alphabetic]{biblatex}
%\usepackage[style=reading]{biblatex}
%\usepackage{csquotes}
%\addbibresource{referencias.bib}

% Definición de entornos teorema, lema, etc.
%\theoremstyle{plain}
\newtheorem{theorem}{Teorema}[chapter]
\newtheorem{lemma}[theorem]{Lema}
\newtheorem{propo}[theorem]{Proposición}
\newtheorem{corollary}[theorem]{Corolario}

\theoremstyle{definition}
\newtheorem{definition}[theorem]{Definición}
\newtheorem{example}[theorem]{Ejemplo}
\theoremstyle{remark}
\newtheorem{remark}[theorem]{Observación}

\let\oldproofname=\proofname
\renewcommand{\proofname}{\rm\bf{Demostración}}
\renewcommand{\qedsymbol}{$\smiley$}


% Configuración de encabezados y pies de página
\pagestyle{fancy}
\fancyhf{}
\fancyhead[LE,RO]{\thepage}  % Número de página en lados opuestos
\fancyhead[RE]{\leftmark}    % Capítulo en páginas de la derecha
\fancyhead[LO]{\rightmark}   % Sección en páginas de la izquierda

% Título y autor del libro
\title{\Huge \textbf{Notas de Geometría Compleja}} 
\author{\Large Héctor Gabriel Rodríguez Morales}
\date{}  % Deja vacío si no quieres que aparezca la fecha en la portada
